% Generated from locale/tr/content/normal-modal-logic/axioms-systems/derived-rules.tex; do not hand-edit.


\olfileid[tr]{nml}{prf}{der}

\olsection{Türetilmiş Kurallar}

!!{derivation}s bulmak ve yazmak açıkça zor, külfetli ve yinelemelidir.
Örneğin çoğu kez $!A \lif !B$'den $\Box !A \lif \Box !B$'ye geçmek, yani
\RK{} kuralını uygulamak isteriz. Bunun için \Nec{} uygulanmalı, sonra
\Ax{K}'nin uygun örneği kaydedilmeli ve ardından \MP{} uygulanmalıdır.
$!A \lif !B$ ile $!B \lif !C$'den $!A \lif !C$'ye geçmek için de (uzun)
totolojik örnek
\[
(!A \lif !B) \lif ((!B \lif !C) \lif (!A \lif !C))
\]
kaydedilip \MP{} iki kez uygulanmalıdır. Çoğu kez bir alt-!!{formula} yerine
ona denk olduğunu bildiğimiz bir formül koymak isteriz; örneğin
\iftag{prvDiamond}{$\Diamond !A$ yerine $\lnot\Box\lnot !A$}{$\Box !A$
yerine $\lnot\Diamond\lnot !A$} ya da $\lnot\lnot !A$ yerine $!A$ koymak
isteriz. Bu yüzden gerçek !!{derivation} metnini yazmak yerine ara adımların neden
!!{derivable} olduğunu kaydetmek daha elverişlidir. Bu amaçla
!!{derivability} hakkında bazı olguları bir araya getirelim.

\begin{prop}
  $\Log{K} \Proves !A_1$, \dots, $\Log{K} \Proves !A_n$ ise ve $!B$,
  $!A_1$, \dots, $!A_n$'den önerme mantığı aracılığıyla çıkıyorsa
  $\Log{K} \Proves !B$ olur.
\end{prop}

\begin{proof}
  $!B$, $!A_1$, \dots,~$!A_n$'den önerme mantığı aracılığıyla çıkıyorsa
  \[
  !A_1 \lif (!A_2 \lif \cdots (!A_n \lif !B)\dots)
  \]
  bir totolojik örnektir. \MP{}'yi $n$ kez uygulamak $!B$ için bir
  !!{derivation} verir.
\end{proof}

Bu önermenin kullanımını \PL{} ile göstereceğiz.

\begin{prop}
  $\Log{K} \Proves !A_1 \lif (!A_2 \lif \cdots (!A_{n-1} \lif
  !A_n)\dots)$ ise $\Log{K} \Proves \Box !A_1 \lif (\Box !A_2 \lif
  \cdots (\Box !A_{n-1} \lif \Box !A_n)\dots)$ olur.
\end{prop}

\begin{proof}
  $n$ üzerinde, \olref[nor]{prop:rk} ispatındaki gibi tümevarımla.
\end{proof}

Bu önermenin kullanımını \RK{} ile göstereceğiz. Şimdi bu sonuçların
!!{derivability} sonuçlarını daha kolay ortaya koymaya nasıl yardım ettiğini
örnekleyelim.

\begin{prop}
  $\Log{K} \Proves (\Box!A \land \Box!B) \lif \Box (!A \land !B)$
\end{prop}

\begin{proof}
  \begin{derivation}
    1. & $\Log{K} \Proves !A \lif (!B \lif (!A \land !B))$ & \Taut \\
    2. & $\Log{K} \Proves \Box!A \lif (\Box !B \lif \Box(!A \land !B)))$ & \RK, 1\\
    3. & $\Log{K} \Proves (\Box!A \land \Box!B) \lif \Box(!A \land !B)$ & \PL, 2
  \end{derivation}
\end{proof}

\begin{prop}\ollabel{prop:rewriting}
  $\Log{K} \Proves !A \liff !B$ ve $\Log{K} \Proves
  \Subst{!C}{!A}{q}$ ise $\Log{K} \Proves \Subst{!C}{!B}{q}$ olur.
  % [tr source emendation] Upstream drops the formula marker in the final
  % substitution argument (`B`); the proposition requires `!B`.
\end{prop}

\begin{proof}
  Alıştırma.
\end{proof}

\begin{prob}
  \olref[nml][prf][der]{prop:rewriting}'i, $!C$'nin karmaşıklığı üzerinde
  tümevarımla, $\Log{K} \Proves !A \liff !B$ ise
  $\Log{K} \Proves \Subst{!C}{!A}{q} \liff \Subst{!C}{!B}{q}$ olduğunu
  ispatlayarak gösterin.
\end{prob}

Bu önerme özellikle $\Diamond$'ı $\Box$'a (ya da tersine) dönüştürmek veya
!!a{formula} içindeki çift değillemeleri kaldırmak istediğimizde kullanışlıdır.
Aşağıda $!C(!A)$ !!a{formula} ifadesini $!C(!B)$ olarak yeniden yazdığımızda
\olref{prop:rewriting} uygulamalarını ``$!A$ yerine $!B$'' ile göstereceğiz.
% [tr source emendation] Upstream calls the displayed C(A)-to-C(B) step
% “A for B,” but its derivation and worked example require “B for A.”
Başka deyişle ``$!A$ yerine $!B$'' şunun kısaltmasıdır:
  \begin{derivation}
    & $\Proves !C(!A)$ \\
    & $\Proves !A \liff !B$\\
    & $\Proves !C(!B)$ & \olref{prop:rewriting} uyarınca
  \end{derivation}
Örneğin:

\begin{prop}
  $\Log{K} \Proves \lnot\Box p \lif \Diamond \lnot p$
\end{prop}

\begin{proof}
  \iftag{prvDiamond}{% Diamond is primitive
  \begin{derivation}
    1. & $\Log{K} \Proves \Diamond \lnot p \liff \lnot\Box\lnot\lnot p$ &
    \Dual\\
    2. & $\Log{K} \Proves \lnot \Box \lnot\lnot p \lif \Diamond \lnot p$ & \PL, 1\\
    3. & $\Log{K} \Proves \lnot\Box p \lif \Diamond \lnot p$ & $\lnot\lnot p$ yerine $p$\\
  \end{derivation}
  }{% Diamond is defined
  \begin{derivation}
    1. & $\Log{K} \Proves \lnot \Box \lnot\lnot p \lif \Diamond \lnot p$ & \Taut\\
    2. & $\Log{K} \Proves \lnot\Box p \lif \Diamond \lnot p$ & $\lnot\lnot p$ yerine $p$\\
  \end{derivation}
  1. satırdaki formül, $q$ yerine $\Box\lnot\lnot p$ koyularak elde edilen
  $\lnot q \lif \lnot q$ örneğidir. $\Diamond$, $\lnot\Box\lnot$ olarak
  tanımlandığından $\Diamond\lnot p$ ile $\lnot\Box\lnot\lnot p$ özdeştir.}
\end{proof}

Yukarıdaki !!{derivation} içinde son adım olan ``$\lnot\lnot p$ yerine $p$'' şu
uzun ifadenin kısaltmasıdır:
  \begin{derivation}
    & $\Log{K} \Proves \lnot \Box \lnot\lnot p \lif \Diamond \lnot
    p$\\
    & $\Log{K} \Proves \lnot\lnot p \liff p$ & \Taut\\
    & $\Log{K} \Proves \lnot\Box p \lif \Diamond \lnot p$ & \olref{prop:rewriting} uyarınca
  \end{derivation}
\olref{prop:rewriting}'deki $!C(q)$, $!A$ ve $!B$ rollerini burada sırasıyla
$\lnot \Box q \lif \Diamond \lnot p$, $\lnot\lnot p$ ve $p$ oynar.

!!a{formula}, $\lnot\Diamond !A$ biçiminde bir alt-!!{formula} içerdiğinde,
$\Log{K} \Proves \lnot\Diamond !A \liff \Box\lnot !A$ olduğundan
\olref{prop:rewriting}'i kullanarak bunun yerine $\Box\lnot !A$ koyabiliriz.
Bunu ve benzeri değiştirmeleri yalnızca ``$\lnot\Diamond$ yerine
$\Box\lnot$'' diye göstereceğiz.

Aşağıdaki önerme, !!{derivability} sonuçlarını şematik olarak ortaya
koyabileceğimizi gerekçelendirir. Örneğin önceki önerme yalnızca
$\Log{K} \Proves \lnot\Box p \lif \Diamond \lnot p$ sonucunu değil, keyfî
$!A$ için $\Log{K} \Proves \lnot\Box !A \lif \Diamond \lnot !A$ sonucunu da
ortaya koyar.

\begin{prop}
  $!A$, $!B$'nin bir yerine koyma örneğiyse ve $\Log{K} \Proves !B$ ise
  $\Log{K} \Proves !A$ olur.
\end{prop}

\begin{proof}
  $!B$ için verilen !!{derivation} uzunluğu üzerinde tümevarımla, bir yerine
  koymayı bütün !!{derivation} boyunca uygulamanın yine doğru bir !!{derivation}
  verdiğini doğrulamak zahmetli fakat sıradandır. Özellikle totolojik
  örneklerin yerine koyma örnekleri yine totolojik örneklerdir;
  \iftag{prvDiamond}{\Dual{} ve~}{}\Ax{K} örneklerinin yerine koyma örnekleri
  yine \iftag{prvDiamond}{\Dual{} ve~}{}\Ax{K} örnekleridir; hem öncül ya da
  öncüllerde hem sonuçta !!{formula}s kullanılarak !!{propositional variable}s
  değiştirildiğinde \MP{} ve \Nec{} uygulamaları doğru kalır.
\end{proof}


